\documentclass[12pt,a4paper]{article}
\usepackage[utf8]{inputenc}
\usepackage[T1]{fontenc}
\usepackage{amsmath, amssymb}
\usepackage{graphicx}
\usepackage{booktabs}
\usepackage{geometry}
\geometry{top=2.5cm, bottom=2.5cm, left=2.5cm, right=2.5cm}

\title{Calibration and Uncertainty Analysis of a Reactive Transport Model in Porous Media}
\author{Your Name \\ \small{University, Department}}
\date{\today}

\begin{document}

\maketitle

\begin{abstract}
This paper presents a numerical model for reactive transport (advection-dispersion with first-order kinetics) and its calibration against synthetic data. Parameter sensitivity and confidence intervals are obtained via Bootstrap. The approach successfully recovers dispersion and reaction coefficients. The code is open-source.
\end{abstract}

\section{Introduction}
Transport of contaminants in groundwater requires knowledge of parameters such as the dispersion coefficient \(D_x\) and reaction rate \(R\). These are often determined by calibrating against field observations.

\section{Mathematical Formulation}
The reactive transport equation:
\begin{equation}
\frac{\partial C}{\partial t} = D_x \frac{\partial^2 C}{\partial x^2} - v \frac{\partial C}{\partial x} - R C,
\end{equation}
with \(C(0,t)=C_{\text{in}}(t)\), \(\partial C/\partial x=0\) at the outlet, and zero initial condition.

\section{Numerical Solution}
A Crank-Nicolson scheme with Thomas algorithm is used. Verification against the Ogata-Banks analytical solution is performed.

\section{Calibration}
The objective function is the root-mean-square error (RMSE) between model and observed outlet concentration. Optimization uses L-BFGS-B. Confidence intervals are obtained via Bootstrap (100 resamples).

\section{Results}
For synthetic data with known parameters (\(D_x = 3\cdot10^{-4}\) m²/s, \(R = 1.5/86400\) s\(^{-1}\)), calibration yields close estimates (Table 1).
\begin{table}[h]
\centering
\caption{Calibration results}
\begin{tabular}{c|c|c}
Parameter & True & Estimated \\
\hline
\(D_x\), m²/s & \(3.0\cdot10^{-4}\) & \(2.98\cdot10^{-4}\) \\
\(R\), s\(^{-1}\) & \(1.74\cdot10^{-5}\) & \(1.71\cdot10^{-5}\) \\
\end{tabular}
\end{table}

95\% confidence intervals: \(D_x\) \([2.7, 3.2]\cdot10^{-4}\), \(R\) \([1.5, 1.9]\cdot10^{-5}\).

\section{Discussion}
Sensitivity analysis shows \(D_x\) mainly controls the spread, while \(R\) influences the tail.

\section{Conclusion}
The developed package enables efficient calibration and uncertainty quantification for reactive transport models and can be applied to real data.

\begin{thebibliography}{9}
\bibitem{ogata} Ogata, A., Banks, R.B. (1961). A solution of the differential equation of longitudinal dispersion in porous media. USGS Professional Paper 411-A.
\bibitem{saltelli} Saltelli, A. et al. (2008). Global Sensitivity Analysis: The Primer.
\end{thebibliography}

\end{document}
